\subsection{\texorpdfstring{PROJECTOR OF \(A^{+}(\mathrm{X})\) ONTO \(\mathrm{L}(\mathrm{X})\)}{PROJECTOR OF A+ (X) ONTO L(X)}}

Let \(\pi\) be the linear mapping of \(A^{+}(X)\) into \(L(X)\) defined by

\[
\pi (x _ {1} \dots x _ {n}) = (\mathrm{ad} (x _ {1}) \circ \dots \circ \mathrm{ad} (x _ {n - 1})) (x _ {n})\tag{3}
\]

for \(n \geqslant 1, x_1, \ldots, x_n\) in X.

PROPOSITION 1. (a) The restriction \(\pi_0\) of \(\pi\) to \(\mathrm{L}(\mathbf{X})\) is a derivation of \(\mathrm{L}(\mathbf{X})\) .

\begin{enumerate}
\def\labelenumi{(\alph{enumi})}
\setcounter{enumi}{1}
\item
  For every integer \(n \geqslant 1\) and all \(u\) in \(\mathbf{L}^n(\mathbf{X})\) , \(\pi(\mathbf{u}) = n.u.\)
\item
  Let E be the endomorphism algebra of the module L(X) and \(\theta\) the homomorphism of A(X) into E such that \(\theta(x) = \operatorname{ad} x\) for all \(x \in X\) . The restriction of \(\theta\) to L(X) is a Lie algebra homomorphism of L(X) into E, which coincides on X with the adjoint representation of L(X), whence
\end{enumerate}

\[
\theta (u). v = [ u, v ] \quad \text { for   } u, v \text {   in   } L (X).\tag{4}
\]

Let \(a\) be in \(\mathrm{A}(\mathrm{X})\) and \(b\) in \(\mathrm{A}^{+}(\mathrm{X})\) ; then

\[
\pi (a. b) = \theta (a). \pi (b).\tag{5}
\]

It suffices to consider the case \(a = x_{1} \ldots x_{p}\) , \(b = x_{p+1} \ldots x_{p+q}\) with \(p \geqslant 0\) , \(q \geqslant 1\) and \(x_{1}, \ldots, x_{p+q}\) in \(\mathbf{X}\) ; but then (5) follows immediately from (3) since \(\theta(x) = \operatorname{ad} x\) for \(x \in \mathbf{X}\) .

Let \(u\) and \(v\) be in \(\mathbf{L}(\mathbf{X})\) ; by (4) and (5),

\[
\begin{array}{r l} \pi_ {0} ([ u, v ]) & = \pi (u v - v u) = \theta (u). \pi (v) - \theta (v). \pi (u) \\ & = [ u, \pi (v) ] - [ v, \pi (u) ] = [ u, \pi_ {0} (v) ] + [ \pi_ {0} (u), v ], \end{array}
\]

hence \(\pi_0\) is a derivation of L(X).

\begin{enumerate}
\def\labelenumi{(\alph{enumi})}
\setcounter{enumi}{1}
\tightlist
\item
  Let \(\pi_1\) be the endomorphism of the module \(\mathrm{L}(\mathbf{X})\) which coincides on \(\mathrm{L}^n (\mathbf{X})\) with multiplication by the integer \(n\geqslant 1\) . The formula \([\mathrm{L}^n (\mathbf{X}),\mathrm{L}^m (\mathbf{X})]\subset \mathrm{L}^{n + m}(\mathrm{X})\) shows that \(\pi_{1}\) is a derivation (Algebra, Chapter III, § 10, no. 3, Example 6). The derivation \(\pi_1 - \pi_0\) of \(\mathrm{L}(\mathbf{X})\) is zero on \(\mathbf{X}\) and, as \(\mathbf{X}\) generates \(\mathrm{L}(\mathbf{X}),\pi_0 = \pi_1\) , whence (b).
\end{enumerate}

COROLLARY. Suppose that K is a Q-algebra. Let P be the linear mapping of \(A^{+}(X)\) into itself such that

\[
\mathrm{P} \left(x _ {1} \dots x _ {n}\right) = \frac {1}{n} (\operatorname{ad} x _ {1} \circ \dots \circ \operatorname{ad} x _ {n - 1}) (x _ {n})\tag{6}
\]

for \(n \geqslant 1\) and \(x_1, \ldots, x_n\) in \(\mathbf{X}\) . Then \(\mathbf{P}\) is a projector of \(\mathbf{A}^+(\mathbf{X})\) onto \(\mathbf{L}(\mathbf{X})\) .

The image of \(\mathbf{P}\) is contained in \(\mathrm{L}(\mathrm{X})\) . Further, for all \(n \geqslant 1\) and all \(u\) in \(\mathrm{L}^n (\mathrm{X})\) , \(\mathrm{P}(u) = \frac{1}{n}\pi (u)\) , whence \(\mathrm{P}(u) = u\) by Proposition 1. As

\[
\mathrm{L} (\mathrm{X}) = \sum_ {n \geqslant 1} \mathrm{L} ^ {n} (\mathrm{X}),
\]

we see that the restriction of P to L(X) is the identity.

Remark. Suppose that K is a field of characteristic zero and let Q be the projector of A(X) = U(L(X)) onto L(X) associated with the canonical graduation of U(L(X)), cf.~§ 1, no. 5. For \(\alpha \in \mathbf{N}^{(\mathbf{X})}\) , Q(A \(^\alpha\) (X)) ⊂ L \(^\alpha\) (X). For it suffices to verify that the image and the kernel of Q are graded submodules of A(X) with the graduation of type N \(^\alpha\) . This is obvious for the image, which is equal to L(X). On the other hand, let n be an integer \(\geqslant\) 1. The vector subspace of A(X) generated by the \(y^n\) , where \(y \in L\) (X), is equal to the vector subspace of A(X) generated by the \(\sum_{\sigma \in \mathcal{E}_n} y_{\sigma(1)} y_{\sigma(2)} \ldots y_{\sigma(n)}\) , where \(y_1, y_2, \ldots, y_n\) are homogeneous elements of L(X); then this subspace is a graded submodule of A(X).

(Note that, if \(\operatorname{Card}(\mathbf{X}) \geqslant 2\) , the projectors \(\mathbf{P}\) and \(\mathbf{Q}\) do not coincide on \(\mathbf{A}^{+}(\mathbf{X})\) . For let \(x, y\) be in \(\mathbf{X}\) with \(x \neq y\) and write

\[
z = x [ x, y ] + [ x, y ] x = x ^ {2} y - y x ^ {2}.
\]

Then \(\mathbf{Q}(z) = 0\) and \(\mathrm{P}(z) = \frac{1}{3}[x, [x,y]] \neq 0\) , cf.~§2, no. 10, Example and no. 11, Theorem 1.)

